% Chapter 5: 논의
\chapter{논의}
\label{ch:discussion}

본 장에서는 제~\ref{ch:results}장의 결과를 해석하고, MutBench의 방법론적 함의, 한계, 향후 연구 방향을 논의한다.
제~\ref{ch:background}장에서 확립된 바와 같이, MutBench는 일반적인 벤치마킹 방법론~\cite{weber2019benchmarking,boulesteix2017towards}의 복합 메트릭 설계와 단일 요인 분리 원칙을 바이러스 유전체학에 적용한다.

%% ============================================================
\section{MutBench의 방법론적 타당성}

%% ============================================================
\subsection{부트스트랩 신뢰 구간의 해석}
\label{sec:bootstrap_interpretation}

영역 수준 BCa 부트스트랩에서 관찰된 넓은 95\% 신뢰 구간(예: MutClust-Hybrid [0.323, 0.635])은 소표본($n = 7$)에서의 영역 의존적 재현율 민감도를 반영한다: 고신호 영역(RBD)이 리샘플에 다수 포함되면 재현율이 상승하고, 저신호 영역만 포함되면 급감한다.
그럼에도 쌍별 부트스트랩 비교에서 방법 간 차이는 유의하며($p < 0.001$), 동일 리샘플에서의 쌍별 평가가 영역 조합 변동을 상쇄하기 때문이다.
향후 ORF1ab, N 단백질 등 추가 도메인을 포함하여 유효 표본 크기를 확대할 필요가 있다.
통계적 견고성이 확립되었으므로, 다음으로 중요한 질문은 합성 벤치마크의 결론이 실제 데이터로 전이되는지 여부이다.

%% ============================================================
\subsection{합성 데이터와 실제 데이터 간의 성능 격차}
\label{sec:synthetic_real_gap}

실제 GISAID 2020--2021 H-score 벤치마크(제~\ref{sec:real_gisaid_benchmark}절)는 합성 결론의 무조건적 외삽이 불가함을 입증하였다.
제~\ref{sec:real_gisaid_benchmark}절에서 제시된 바와 같이, 실제 H-score의 희소 분포(29,903개 중 342개 비영)는 탐지 역학을 근본적으로 변화시킨다.
밀도 기반 클러스터링은 이러한 극단적 희소 분포에서 클러스터 형성이 제한되며(HDBSCAN-Auto는 클러스터 미형성), 합성 벤치마크에서의 MutClust-Hybrid 대비 우위(hotspot-score 0.778 vs. 0.646)도 희소 분포에서는 사라져 두 방법의 성능이 수렴한다(0.383 vs. 0.386).

그럼에도, 실제 GISAID 데이터에서도 농축비가 무작위 기준선 대비 6.9--8.0배를 유지하여, 상위 방법이 포착하는 생물학적 신호는 MCC/F1의 실질적 하락에도 불구하고 보존됨을 보여준다.
합성 벤치마크는 통제된 신호 복원 시험으로 설계되었기에 그 가치를 유지한다: 목적은 실제 절대 성능을 예측하는 것이 아니라, 교란 요인(희소성, 창시자 효과, 표본 편향)이 통제된 조건에서 상대적 방법 강점을 드러내는 것이다.
나아가, 합성과 실제를 함께 평가하는 이중 벤치마크 설계 자체가 방법론적 기여이다: 이론적 탐지 역량과 실용적 성능 간의 격차를 정량화하며, 이는 어느 한쪽 벤치마크만으로는 얻을 수 없는 정보이다.
따라서 MutBench는 방법의 ``할 수 있는 것''(합성)과 ``실제로 하는 것''(실제)을 모두 평가하되, 측정된 격차를 한계가 아닌 보정 계수로 취급할 것을 권장한다.

%% ============================================================
\subsection{Ground Truth와 탐지 방법의 독립성}
\label{sec:circularity}

합성 벤치마크의 ground truth는 탐지 방법과 독립적으로 정의되었다(DMS 실험 데이터 및 수렴 진화 문헌). 이는 평가와 탐지 간의 순환성을 배제한다.
Min-max 정규화를 MutBench 전체에서 기본값으로 사용하였으며, z-score 및 순위 변환은 희소 데이터 분포에서 인공물을 유발할 수 있어 평가하지 않았다.
탐지된 핫스팟의 생물학적 타당성은 3D 구조 분석(제~\ref{sec:structure_analysis}절)에서 5.92배 공간 접촉 농축이 확인됨으로써 추가로 뒷받침되었으며, 제~\ref{sec:external_comparison}절에서 도메인 특이적 방법(MutClust 변이체)이 범용 대안(OPTICS, KDE)보다 일관되게 우수함이 확인되었다.
벤치마크 타당성이 확립되었으므로, 기반이 되는 H-score 지표의 시간적 한계를 검토할 필요가 있다.

%% ============================================================
\section{H-score 기반 탐지의 한계}

%% ============================================================
\subsection{고정 참조 기반 점수의 한계}
\label{sec:hscore_saturation}

시간적 검증(제~\ref{sec:temporal}절)에서 관찰된 H-score 포화 현상은 참조 기반 지표의 근본적 한계를 드러낸다: 2024--2025 데이터에서 모든 1,273개 Spike 위치가 비영 H-score를 가지며 표준 편차는 0.098에 불과하여 판별력을 상실한다.
TC-freq가 수렴 진화 ground truth에서 최고 농축비(3.43배)를 달성한 것은 시간적 대비 점수의 타당성을 시사하나, 최적 시간 해상도의 결정은 향후 과제로 남는다.
포화 현상은 다중 독립 ground truth 정의에 대한 평가의 동기를 제공한다.

%% ============================================================
\subsection{상보적 평가 관점}
\label{sec:complementary_gt}

다중 ground truth 평가(제~\ref{sec:multi_gt}절)에서 DMS-escape와 수렴 진화 위치 간 Jaccard 유사도가 0.006에 불과하여, 두 측정이 변이의 본질적으로 다른 측면을 포착함을 확인하였다.
단일 ground truth만 사용하면 특정 생물학적 차원에 부합하는 방법이 과대평가될 위험이 있으며, MutBench의 다중 ground truth 접근법은 이를 완화한다.
최우수 방법이 ground truth에 따라 달라진다는 결과는 ``최고의 핫스팟 탐지 방법''이 맥락 의존적임을 시사한다: 면역 회피 연구에서는 수렴 진화 기준, 분자 기능 연구에서는 DMS 기준에 최적인 방법이 각각 선호된다.
MutBench는 이러한 맥락 의존적 방법 선택을 체계적으로 안내한다.
SARS-CoV-2 내에서 관찰된 ground truth 의존적 방법 순위는 9-병원체 벤치마크에서 드러난 병원체 의존적 순위를 예고한다.

%% ============================================================
\section{교차 병원체 일반화와 PAHD}

%% ============================================================
\subsection{교차 병원체 일반화의 함의}
\label{sec:cross_pathogen_implications}

제~\ref{ch:results}장에서 상술한 바와 같이, 9-병원체 벤치마크는 범용 최적 방법이 존재하지 않으며 모든 병원체가 고유한 최적 조합을 보이고 상호작용 효과가 상당함을 입증하였다. 이 결과는 Nextstrain~\cite{nextstrain2018hadfield}과 같은 실시간 병원체 감시 시스템의 맥락에서 중요한 함의를 가진다: 새로운 병원체 또는 변이체 출현 시 기존 최적 조합을 그대로 적용하는 것은 부적절하며, 병원체 특성에 기반한 적응적 방법 선택이 필요하다.
근본 원인은 병원체별 게놈 특성---엔트로피 분포, 보존 영역 비율, 양성 선택 부위 밀도---이 탐지 방법에 대해 근본적으로 상이한 신호 지형을 형성하기 때문이다.
예를 들어, 전체적으로 높은 변이율을 가진 병원체(예: H3N2, 모든 위치가 비영 H-score)는 역치 기반 시드 탐지의 이점을 무력화하는 반면, 컴팩트한 게놈과 소수의 수렴 진화 부위를 가진 병원체(예: MERS)는 희소 데이터에서 약한 신호를 탐지할 수 있는 방법을 선호한다.
개념 증명 시뮬레이션은 빈도 기반 스코어링이 창시자 효과와 양성 선택을 혼동함을 추가로 확인하였으며($\rho = -0.876$), 병원체 의존적 차원을 하나 더 추가한다: 창시자 효과 오염 정도가 각 병원체의 전파 이력에 따라 달라진다.
이러한 결과에 기반하여, PAHD(Pathogen-Adaptive Hotspot Detection)가 병원체별 방법 선택을 위한 개념적 프레임워크로 제안되며, Rice의 알고리즘 선택 문제~\cite{rice1976algorithm}, 즉 자동화된 알고리즘 선택~\cite{kerschke2019algorithm} 및 AutoML~\cite{kotthoff2016algorithm}로 확장되어 온 프레임워크로 형식화될 수 있다.
구체적으로, MutBench의 결과는 Rice 프레임워크를 다음과 같이 실체화한다: \textit{문제 공간} $\mathcal{P}$는 게노믹 프로파일(정렬 길이, 다양성 비율, ground truth 밀도 등)로 특성화된 병원체로 구성되고, \textit{알고리즘 공간} $\mathcal{A}$는 351개의 scoring--detection 조합(9 $\times$ 39)으로 구성되며, \textit{성능 공간} $\mathcal{Y}$는 3-layer ground truth에 대해 평가된 MCC이다. 상호작용 우세($\omega^2 = 0.285$)는 매핑 $f: \mathcal{P} \to \mathcal{A}$가 비자명함을 실증적으로 확인하며, 이는 어떤 단일 알고리즘도 모든 가능한 문제 분포에서 다른 모든 알고리즘을 능가할 수 없다는 No Free Lunch 정리~\cite{wolpert1997nfl}의 직관과 일치하는 발견이다. PAHD 프로토타입의 9개 중 5개 양의 MCC 결과는 프로파일 공간 $\mathcal{P}$에서 단순한 최근접 이웃 접근법만으로도 이 매핑을 부분적으로 복원할 수 있음을 보여주며, 더 정교한 메타 학습 접근법의 동기를 제공한다.

PAHD 프로토타입 결과(제~\ref{subsec:pahd_prototype}절)는 더 거친 예측 단위에서 oracle 성능의 최대 57.7\%를 회복할 수 있음을 보여주나, $n = 9$의 한계는 여전히 존재한다.
ESM-2 교차 패러다임 검증(제~\ref{subsec:esm_validation}절)은 이 발견을 더욱 강화한다: ESM-LLR은 SARS-CoV-2와 MERS에서 1위를 차지하였으나 Norovirus와 HIV-1에서는 최하위를 기록하여, MSA 기반이든 단백질 언어 모델 기반이든 어떤 단일 스코어링 패러다임도 모든 병원체에서 우위를 점하지 못함을 확인하였다.
PAHD의 구체적 구현은 제~\ref{ch:conclusion}장의 향후 연구로 제안한다.
입증된 병원체 적응적 선택의 필요성은 달성된 절대 성능 수준의 해석이라는 과제를 제기한다.

%% ============================================================
\subsection{탐지 성능의 상한}
\label{sec:oracle_mcc_ceiling}

9개 병원체에 걸친 평균 oracle MCC는 0.298(범위: 0.088--0.549)이며, 이 값이 낮아 보일 수 있으나 세 가지 구조적 요인으로 설명된다.
첫째, \textbf{위치 수준 정확 매칭}은 단 한 위치만 벗어난 탐지도 위양성으로 처리하는 엄격한 기준을 부과한다.
둘째, \textbf{극단적 클래스 불균형}(양성 비율 대부분 10\% 미만)이 달성 가능한 MCC를 제약한다---양성 비율 5\%에서 90\% 재현율과 30\% 정밀도는 이론적 MCC 0.49만을 산출한다.
셋째, \textbf{ground truth의 엄격성}으로 인해 실제이지만 미주석된 핫스팟을 탐지하는 방법이 불이익을 받는다. 문헌에서 확인된 수렴 진화 부위만이 양성으로 처리되기 때문이다.

이러한 제약에도 불구하고, oracle MCC 0.298은 명백히 비무작위이다: 무작위 기준선(10,000회 순열)의 평균은 0.001(SD = 0.003)이며, oracle은 무작위 기대보다 약 99 표준편차 위에 위치한다.
맥락적으로, 다른 생물정보학 벤치마크에서의 유사한 위치 수준 평가 지표도 비슷한 범위를 보인다: 전사인자 결합 부위 예측은 단일 뉴클레오티드 해상도에서 F1 0.3--0.5를 달성하며~\cite{weirauch2013evaluation}, 초기 CASP 단백질 구조 예측은 난이도 높은 표적에서 GDT-TS 0.3 미만을 기록하였다~\cite{moult2014casp}(GDT-TS는 이진 분류가 아닌 구조적 유사성을 측정하나, 위치 수준 생물학적 예측에서 낮은 절대 점수가 일반적임을 보여주는 비교이다).
따라서 oracle MCC 0.298은 위치 수준 생물학적 예측 과제의 본질적 난이도와 일치한다.

위치 정확 평가의 영향을 정량화하기 위해, $\pm5$ 및 $\pm10$ 위치 윈도우를 사용한 영역 중첩 분석을 수행하였다(표~\ref{tab:region_overlap_mcc}).
영역 중첩 평가에서는 탐지된 위치가 지정된 윈도우 내에 ground truth 위치를 포함하면 참양성으로 계산되고, ground truth 위치는 윈도우 내에 탐지가 존재하면 재현된 것으로 계산된다.
$\pm5$ 윈도우는 대략 단일 단백질 이차 구조 요소 또는 항원 에피토프({$\sim$}15개 잔기)에 해당하여 생물학적으로 의미 있는 근접성 기준을 제공한다. 이 분석은 위치 정확 평가의 엄격성에 대한 민감도 검사로 기능하며, 대체가 아니다: 위치 정확 MCC는 매개변수가 없고 완전한 재현성을 보장하므로 주요 지표로 유지된다.
$\pm10$ 윈도우 평가에서 평균 oracle MCC는 0.289에서 0.712로 146\% 향상되었다.
특히 H3N2는 $\pm10$에서 MCC = 0.964(정확 매칭 0.521 대비)를 달성하여, 탐지된 핫스팟이 정확히 중첩되지 않더라도 알려진 항원 변이 부위에 공간적으로 근접함을 확인하였다.
이는 위치 정확 평가의 보수성을 정량화한다: 완화된 윈도우에서의 대폭적인 향상은 많은 탐지가 게놈 전체에 무작위로 분산되지 않고 ground truth 위치의 생물학적으로 관련 있는 근접 범위 내에 위치함을 나타낸다.

\begin{table}[htbp]
\centering
\caption{세 가지 평가 모드에서의 oracle 조합 영역 중첩 MCC 분석}
\label{tab:region_overlap_mcc}
\small
\begin{tabular}{lccc}
\toprule
\textbf{병원체} & \textbf{정확 매칭} & \textbf{$\pm$5 윈도우} & \textbf{$\pm$10 윈도우} \\
\midrule
SARS-CoV-2    & 0.227 & 0.712 & 0.861 \\
H3N2          & 0.521 & 0.851 & 0.964 \\
Norovirus     & 0.355 & 0.925 & 0.925 \\
HIV-1         & 0.301 & 0.736 & 0.789 \\
Dengue        & 0.274 & 0.741 & 0.919 \\
RSV           & 0.150 & 0.164 & 0.164 \\
Influenza B   & 0.195 & 0.558 & 0.594 \\
MERS          & 0.071 & 0.263 & 0.295 \\
HCV           & 0.510 & 0.793 & 0.894 \\
\midrule
\textbf{평균} & \textbf{0.289} & \textbf{0.638} & \textbf{0.712} \\
\bottomrule
\end{tabular}

\vspace{1mm}
{\footnotesize $^\dagger$영역 중첩 분석의 위치 정확 MCC(0.289)는 벤치마크 보고 oracle(0.298)과 약간 차이가 나는데, 이는 영역 중첩 분석이 경계 처리(예: 정렬 경계의 가장자리 위치 처리)에서 소규모 차이가 있는 재구현된 평가 파이프라인을 사용하기 때문이다. 절대적 기준선 차이(0.009, 즉 3.1\%)는 주요 발견에 영향을 미치지 않는다: 위치 정확에서 영역 중첩으로의 상대적 향상(146\%)은 탐지된 핫스팟이 ground truth 위치에 공간적으로 근접함을 보여준다.}
\end{table}

\begin{figure}[htbp]
\centering
\includegraphics[width=\textwidth]{region_overlap_bar.png}
\caption{각 병원체에 대한 세 가지 평가 모드(정확 매칭, $\pm$5, $\pm$10 윈도우)에서의 영역 중첩 MCC. 정확 매칭에서 $\pm$10 윈도우 평가로의 일관된 향상은 탐지된 핫스팟이 9개 병원체 전체에서 ground truth 위치에 공간적으로 근접함을 보여준다.}
\label{fig:region_overlap_bar}
\end{figure}

Ground truth가 잘 특성화된 병원체에서는 실질적으로 높은 값을 달성한다.
이러한 상한 효과에도 불구하고, 벤치마크 결과는 공중보건 감시에 대한 실용적 함의를 가진다.

%% ============================================================
\section{실용적 함의와 근본적 한계}

제~\ref{ch:introduction}장에서 제기한 세 가지 문제로 돌아가면: 표준화된 평가 프레임워크의 부재는 MutBench의 3계층 ground truth와 hotspot-score 지표로 해결되었으며(기여 1); 교차 병원체 일반화 근거의 부족은 네 가지 독립적인 통계적 증거에 의해 해소되었고(기여 2); 탐지 방법의 병원체 의존성은 ANOVA 상호작용 우세($\omega^2 = 0.285$)와 LOPO 실패(0/9)로 정량화되어, PAHD 프레임워크의 동기를 제공하였다.

%% ============================================================
\subsection{공중보건 감시에서의 실용적 고려}
\label{sec:practical_surveillance}

SARS-CoV-2의 세대 시간(약 5일)과 서열 제출 지연(중앙값 약 14일)을 감안하면, 실시간 TC-freq 경보에는 최소 4주 윈도우가 필요하다.
조기경보에는 HDBSCAN-Auto(높은 재현율), 백신 후보 선정에는 MutClust-Hybrid(균형 성능)가 적합하며, 회고적 분석에서 MutClust-Hybrid의 높은 재현율은 신흥 핫스팟 위치의 조기 탐지 가능성을 시사하나, 정량적 시점 추정에는 별도의 시간적 검증이 필요하다.

2단계 탐지 파이프라인---HDBSCAN-Auto 1차 선별 후 MutClust-Hybrid 2차 확정(절~\ref{sec:variant_comparison})---은 정밀도(0.968)를 유지하면서 안정성이 0.740으로 향상되어 최고 hotspot-score(0.784)를 달성하였다.

\begin{table}[htbp]
\centering
\caption{2단계 파이프라인 검증 결과 (합성 데이터, 29,903 위치, 1,251개 ground truth 위치)}
\label{tab:twostep_pipeline}
\small
\begin{tabular}{lrrrrrr}
\toprule
\textbf{접근법} & \textbf{정밀도} & \textbf{재현율} & \textbf{F1} & \textbf{탐지 수} & \textbf{안정성} & \textbf{HS} \\
\midrule
HDBSCAN-Auto (단독) & 0.079 & 0.944 & 0.146 & 14,916 & 0.788 & 0.604 \\
MutClust-Hybrid (단독) & 0.968 & 0.659 & 0.785 & 852 & 0.706 & 0.778 \\
\midrule
\textbf{2단계 교집합} & \textbf{0.968} & 0.643 & 0.772 & 831 & \textbf{0.740} & \textbf{0.784} \\
\bottomrule
\end{tabular}
\end{table}

백신 설계 등 후속 실험 비용이 높은 응용에서는 이 2단계 접근법이, 일상 감시에는 MutClust-Hybrid 단독이, 조기경보에는 HDBSCAN-Auto 단독이 각각 적합하다.
다만, 실제 GISAID 2020--2021 데이터에서는 HDBSCAN-Auto가 비영 위치의 희소성(342개)으로 인해 클러스터를 형성하지 못하여, 2단계 파이프라인의 이점이 발현되지 않았다.
모든 빈도 기반 접근법에 영향을 미치는 근본적인 방법론적 한계는 계통발생 비독립성이다.

%% ============================================================
\subsection{계통발생 비독립성과 확장성 한계}
\label{sec:phylo_limits}

본 연구의 모든 탐지 방법은 서열 간 독립성을 암묵적으로 가정하나, SARS-CoV-2 서열은 계통발생적으로 구조화되어 있어 빈도 기반 접근법만으로는 양성 선택과 founder effect를 구분할 수 없다(예: D614G;~\cite{korber2020tracking}).
dN/dS 필터는 이를 부분적으로 해소하나, dN/dS 자체도 계통 구조에 의해 편향될 수 있다.

개념 증명 시뮬레이션(제~\ref{subsec:phylo_correction_poc}절)에서 Fitch 절약법 기반 homoplasy 점수가 수렴 변이를 정확히 최상위로 배치하였으며, TreeTime ML 확장 검증(제~\ref{subsec:treetime_ml_validation}절)에서 동등한 정확도(AUC = 1.000)를 확인하였다.
실제 GISAID 서열에 대한 경험적 검증(제~\ref{subsec:treetime_gisaid_validation}절)에서 H-score와 homoplasy 점수가 근본적으로 다른 신호를 포착함을 확인하여($\rho = -0.107$, $p = 1.36 \times 10^{-4}$), homoplasy 기반 보정의 실용적 가치를 뒷받침하였다.
Homoplasy 민감도 분석(제~\ref{subsec:homoplasy_sensitivity}절)은 2020-H1 서열 데이터와 Layer~A 수렴 진화 ground truth 간의 시간적 불일치를 드러내어, 계통발생 보정에는 시간적으로 정렬된 데이터가 필요함을 확인하였다.
완전한 계통발생 보정과 시간적으로 적절한 서열 표본 추출의 파이프라인 통합은 향후 핵심 확장 과제로 남아 있다(제~\ref{ch:conclusion}장).

%% ============================================================
\subsection{GISAID 데이터의 대표성과 표본 편향}

GISAID~\cite{shu2017gisaid} 기반 입력 데이터에는 체계적 편향이 존재한다.
\textbf{지리적 편향}: 고소득 국가가 전체 서열의 80\% 이상을 기여하여~\cite{chen2022gisaid}, 아프리카/남아시아의 고유 변이가 과소 반영되며, 주별 층화는 시간적 불균형만 보정한다.
\textbf{시간적 편향}: 초기 팬데믹과 Omicron 시기의 서열 축적량 차이가 크며, 주 내 대표성은 보장되지 않는다.
\textbf{Founder effect}: JN.1 확산에 동반된 ``hitchhiker'' 변이~\cite{kaku2024virological}처럼, 기능적 이점 없이 고빈도에 도달하는 위치를 양성 선택과 구분하는 것은 빈도 기반 접근법의 근본적 한계이다.

%% ============================================================
\subsection{범위 정당화: 9개 RNA 바이러스}

9종 병원체는 6개 주요 RNA 바이러스 과(Coronaviridae, Orthomyxoviridae, Retroviridae, Pneumoviridae, Flaviviridae, Caliciviridae)를 체계적으로 포괄하도록 선정되었으며, 팬데믹 잠재력을 가진 RNA 바이러스의 주요 분류학적 다양성을 대표한다.
RNA 바이러스로의 범위 제한은 의도적인 설계 결정이다: H-score 기반 스코어링은 충분한 변이 밀도를 요구하며, RNA 바이러스는 높은 변이율($10^{-3}$--$10^{-5}$ substitutions/site/year)로 이를 제공하는 반면, DNA 바이러스($10^{-6}$--$10^{-8}$)는 근본적으로 다른 스코어링 접근법을 필요로 한다.
표본 크기($n=9$)의 제한에도 불구하고, Friedman 검정은 $p=0.0015$를 달성하여, 효과 크기가 제한된 수의 병원체로도 탐지 가능할 만큼 충분히 크다는 것을 보여준다.

%% ============================================================
\subsection{ESM-2 학습 데이터 중복}
\label{sec:esm2_leakage}

ESM-2~\cite{lin2023esm2}는 바이러스 단백질 서열을 포함하는 UniRef50/UniRef90으로 학습되었으며, 본 벤치마크의 9개 대상 단백질(예: SARS-CoV-2 Spike, HIV-1 gp120, H3N2 HA)은 학습 코퍼스에 포함되어 있을 가능성이 높다.
따라서 ESM-2 log-likelihood ratio(LLR) 점수는 순수한 일반화 가능한 단백질 언어 이해가 아닌 기억된 진화적 패턴을 부분적으로 반영할 수 있으며, 이는 절대 AUC 값을 상향 편향시킬 수 있다.
그러나 이 우려는 분석의 9개 병원체 모두에 동일하게 적용되므로, 절대 성능 수준은 편향될 수 있으나 본 연구의 주요 발견인 교차 병원체 간 상대적 비교의 타당성은 유지된다.
의미 변화 특성(야생형과 돌연변이 임베딩 간 코사인 거리)은 절대 우도가 아닌 상대적 임베딩 변화를 측정하므로 학습 데이터 기억 효과에 대해 더 강건하다.
향후 바이러스 서열을 제외한 데이터베이스로 학습된 단백질 언어 모델을 평가하여 이 잠재적 편향의 크기를 정량화할 필요가 있다.

%% ============================================================
\subsection{병원체 간 Ground Truth 이질성}
\label{sec:gt_heterogeneity}

3계층 ground truth 프레임워크의 핵심 한계는 Layer~A(양성 선택 부위)가 병원체별로 정의되는 방식의 이질성이다.
SARS-CoV-2와 H3N2의 경우, Layer~A 위치는 계통발생적으로 확인된 수렴 진화 부위---동일한 돌연변이가 독립적으로 여러 계통에서 발생한 위치---이다.
반면, HCV의 Layer~A는 수렴적이 아닌 다양화 선택을 보이는 초가변 영역(HVR1/HVR2)으로 구성되며, Norovirus의 Layer~A는 항원 전환 사건과 관련된 시대적 변이 위치로 구성되고, 나머지 병원체는 다양한 수준의 계통발생적 확인을 거친 문헌 보고 면역 회피 부위를 사용한다(표~\ref{tab:gt_positions}).
이러한 이질성은 ANOVA에서 방법 $\times$ 병원체 상호작용을 인위적으로 부풀릴 수 있다. 병원체 생물학의 차이뿐만 아니라 ground truth 정의 기준의 차이도 관찰된 상호작용 효과($\omega^2 = 0.285$)에 기여할 수 있기 때문이다.
그러나 이 이질성은 핫스팟 탐지의 현실적 도전을 반영한다: ground truth 정의는 필연적으로 병원체별 생물학에 의존하며, 단일 기준(예: 수렴 진화만)으로는 모든 병원체에 적용할 수 없다.
탐지 성능이 병원체에 의존한다는 발견은 이 의존성이 병원체 간 본질적 생물학적 차이에서 비롯되든, ground truth 구성의 차이에서 비롯되든, 또는 양자의 상호작용에서 비롯되든 관계없이 유효하다---모든 경우에 실용적 함의는 동일하다: 단일 탐지 방법을 병원체 전체에 균일하게 적용하면 성능 저하가 발생한다.
또한, HCV는 유효 Layer~B 크기가 0이다(보존 영역이 정렬 길이 내에서 매핑될 수 없었음). 이는 HCV 평가가 Layer~A에만 의존함을 의미하며, HCV 특이적 결과를 해석할 때 이 비대칭성을 고려해야 한다.
